% PlumeWatch manuscript skeleton
%
% Working title:
% PlumeWatch: Physics-Guided Deep Learning for Detecting NOx Emissions Changes
% from Power Plants Using TEMPO Satellite Imagery

\documentclass[11pt]{article}

% -----------------------------------------------------------------------------
% Manuscript formatting
% -----------------------------------------------------------------------------
\usepackage[letterpaper,margin=1in]{geometry}
\usepackage[T1]{fontenc}
\usepackage[utf8]{inputenc}
\usepackage{lmodern}
\usepackage{microtype}
\usepackage{setspace}
\onehalfspacing

% Mathematics, tables, and figures
\usepackage{amsmath}
\usepackage{amssymb}
\usepackage{booktabs}
\usepackage{graphicx}
\usepackage{xcolor}

% Citations and links
\usepackage[round,authoryear]{natbib}
\usepackage[hidelinks]{hyperref}
\usepackage[nameinlink,capitalise]{cleveref}

% Keep figures and tables near their first reference during drafting.
\usepackage{placeins}

% -----------------------------------------------------------------------------
% Draft helpers
% -----------------------------------------------------------------------------
\newcommand{\todo}[1]{\textcolor{red}{[TODO: #1]}}
\newcommand{\NOx}{NO$_x$}
\newcommand{\modelname}{PlumeWatch}

\newcommand{\figureplaceholder}[2]{%
  \begin{figure}[t]
    \centering
    \fbox{%
      \parbox[c][1.6in][c]{0.88\linewidth}{%
        \centering \texttt{#1}%
      }%
    }
    \caption{#2}
    \label{fig:#1}
  \end{figure}%
}

% -----------------------------------------------------------------------------
% Title block
% -----------------------------------------------------------------------------
\title{\modelname: Physics-Guided Deep Learning for Detecting
\texorpdfstring{\NOx}{NOx} Emissions Changes from Power Plants Using TEMPO
Satellite Imagery}

\author{%
  Author Name(s)\\
  Affiliation(s)\\
  \texttt{corresponding.author@example.edu}
}

\date{Draft: \today}

\begin{document}

\maketitle

\begin{abstract}
\todo{Write a 150--250 word abstract covering the problem, data, method,
evaluation design, principal result, and practical interpretation.}
\end{abstract}

\noindent\textbf{Keywords:}
TEMPO; nitrogen oxides; power plants; emissions-change detection;
physics-guided machine learning; remote sensing

% -----------------------------------------------------------------------------
\section{Introduction}
\label{sec:introduction}

% TODO: Motivate near-hourly monitoring of point-source emissions.
% TODO: State why column NO2 imagery alone makes change detection difficult.
% TODO: Define the three-class task: decrease, steady, and increase.
% TODO: Summarize the proposed model and evaluation design.
% TODO: End with a concise, numbered contribution list.

\section{Related Work}
\label{sec:related-work}

\subsection{Satellite-Based Monitoring of Point-Source Air Pollution}
% TODO: Cover NO2 retrievals, plume attribution, and geostationary observations.

\subsection{Deep Learning for Atmospheric Remote Sensing}
% TODO: Cover spatial and spatiotemporal models for satellite retrieval analysis.

\subsection{Physics-Guided Learning for Atmospheric Transport}
% TODO: Position explicit advection residuals relative to purely learned models.

\section{Data and Problem Formulation}
\label{sec:data}

\subsection{Data Sources}
% TODO: Describe TEMPO NO2, HRRR weather, and CAMPD emissions records.

\begin{table}[t]
  \centering
  \caption{Data sources and their roles in the study.}
  \label{tab:data-sources}
  \begin{tabular}{lll}
    \toprule
    Source & Variables & Role \\
    \midrule
    TEMPO & \todo{Complete} & Model input \\
    HRRR & \todo{Complete} & Weather and transport \\
    CAMPD & \todo{Complete} & Labels and operations events \\
    \bottomrule
  \end{tabular}
\end{table}

\subsection{Areas of Interest and Sample Construction}
% TODO: Define the 72 km AOIs, five-scan windows, and raster resolution.

\subsection{Emissions-Change Labels}
% TODO: Define early and late time-weighted NOx means and the class threshold.
% TODO: Include the label equation and units.

\subsection{Geographically Disjoint Splits}
% TODO: Describe overlap clustering, split proportions, and frozen test usage.
% TODO: Report record and AOI counts by split and class.

\subsection{Plant-Operations Evaluation Cohort}
% TODO: Define start and stop events, matching rules, cohort size, and purpose.
% TODO: State that this pooled validation/test cohort is a diagnostic evaluation.

\section{Methods}
\label{sec:methods}

\figureplaceholder{architecture}{Model overview. Replace this box with the
final architecture diagram showing observation, transport, context, temporal,
classification, and auxiliary-raster paths.}

\subsection{Raster Completion and Quality Encoding}
% TODO: Describe masked-model filling and the retained NO2 quality channel.

\subsection{Seasonal and Solar Baseline}
% TODO: Define day-of-year and local-solar-hour features and the baseline MLP.

\subsection{Physics-Guided Transport Residuals}
% TODO: Define backward semi-Lagrangian advection and residual quality.
% TODO: Document wind blending, dispersion, decay, and physical units.

\subsection{Two-Branch Spatiotemporal Encoder}
% TODO: Describe observation and residual encoders, aligned fusion, and ConvGRU.

\subsection{Seasonal Context Modulation}
% TODO: Define identity-initialized FiLM scale and shift at each timestep.
% TODO: Describe context dropout and shuffled-context evaluation.

\subsection{Classification and Auxiliary Raster Objectives}
% TODO: Define weighted cross-entropy, raster L1 loss, and total objective.
% TODO: Describe target construction, normalization, and smoothing.

\section{Experimental Design}
\label{sec:experiments}

\subsection{Training and Model Selection}
% TODO: Report optimizer, learning rates, batch size, seed, and early stopping.
% TODO: State that validation macro F1 selects checkpoints.

\subsection{Comparisons and Ablations}
% TODO: List the seasonal baseline and context-aware vision model.
% TODO: List architecture, transport, context, input, and loss ablations.

\subsection{Evaluation Metrics}
% TODO: Define macro F1, balanced accuracy, per-class recall, confusion matrices,
% predicted-class frequencies, and any confidence interval procedure.

\subsection{Reproducibility Controls}
% TODO: Document training-only normalization, fixed seeds, paired records,
% software versions, hardware, and run identifiers.

\section{Results}
\label{sec:results}

\subsection{Primary Held-Out Test Results}
% TODO: Report seasonal and vision results on the frozen test split.

\begin{table}[t]
  \centering
  \caption{Primary held-out test performance. Replace placeholders after the
  final model run.}
  \label{tab:primary-results}
  \begin{tabular}{lccccc}
    \toprule
    Model & Macro F1 & Balanced accuracy & Decrease F1 & Steady F1 & Increase F1 \\
    \midrule
    Seasonal baseline & \todo{TBD} & \todo{TBD} & \todo{TBD} & \todo{TBD} & \todo{TBD} \\
    \modelname{} & \todo{TBD} & \todo{TBD} & \todo{TBD} & \todo{TBD} & \todo{TBD} \\
    \bottomrule
  \end{tabular}
\end{table}

\subsection{Ablation Results}
% TODO: Report controlled comparisons on identical records and fixed seeds.

\subsection{Seasonal-Shortcut Diagnostics}
% TODO: Report context dropout and shuffled-context results.

\subsection{Performance by AOI and Observation Conditions}
% TODO: Stratify results by plant characteristics, month, solar hour, and quality.

\subsection{Plant-Operations Evaluation}
% TODO: Report per-class F1 for seasonal and vision models.
% TODO: Discuss cohort composition and AOI-clustered uncertainty.

\figureplaceholder{plant-operations-f1}{Seasonal and vision per-class F1 on the
plant-operations evaluation cohort.}

\section{Discussion}
\label{sec:discussion}

% TODO: Interpret where imagery adds information beyond seasonal context.
% TODO: Discuss transport assumptions and errors on operations-linked records.
% TODO: Separate supported conclusions from hypotheses for future experiments.

\section{Limitations}
\label{sec:limitations}

% TODO: Cover label timing, retrieval uncertainty, masked filling, transport
% simplifications, geographic scope, class thresholds, and cohort selection.

\section{Conclusion}
\label{sec:conclusion}

% TODO: Summarize the task, method, principal finding, and next validation step.

\section*{Data and Code Availability}
% TODO: Add data access conditions, repository URL, license, and archived version.

\section*{Acknowledgments}
% TODO: Add funding, computing resources, data providers, and contributors.

\section*{Author Contributions}
% TODO: Add a CRediT contribution statement.

\section*{Competing Interests}
% TODO: Add the required disclosure statement.

% Add docs/references.bib, then uncomment these lines.
% \bibliographystyle{plainnat}
% \bibliography{references}

\appendix

\section{Supplementary Dataset Details}
% TODO: Add quality-control rules, split tables, and label sensitivity checks.

\section{Supplementary Model and Training Details}
% TODO: Add layer dimensions, hyperparameters, transport settings, and compute.

\section{Supplementary Results}
% TODO: Add full confusion matrices, predicted-class frequencies, and per-AOI results.

\end{document}
